\documentclass[11pt,a4paper]{article}
\usepackage[utf8]{inputenc}
\usepackage[T1]{fontenc}
\usepackage{geometry}
\geometry{left=2.7cm,right=2.7cm,top=2.6cm,bottom=2.6cm}
\usepackage{lmodern,microtype,setspace}
\usepackage{graphicx,float,amsmath,amssymb,cite,listings,xcolor,booktabs,array,caption,subcaption,algorithm,algpseudocode,tikz,pgfplots,seqsplit}
\usepackage[hidelinks]{hyperref}
\pgfplotsset{compat=1.18}
\usetikzlibrary{shapes,arrows,positioning,fit,backgrounds,matrix,calc}
\onehalfspacing
\setlength{\parindent}{1.5em}
\setlength{\parskip}{0.15em}
\hypersetup{
  pdftitle={Privacy Lens: An Evidence-Bounded Android Privacy Review Framework},
  pdfauthor={Zhiheng Zhang}
}
\title{Privacy Lens: An Evidence-Bounded Android Privacy Review Framework with Replaceable Regulation Packs (Revision 10)}
\author{Zhiheng Zhang}
\date{August 2026}
\begin{document}
\maketitle
\begin{abstract}
This thesis presents Privacy Lens, a local-first Android prototype that converts bounded permission-audit summaries into explainable prompts for human privacy review. The research addresses a recurring problem in mobile privacy tooling: technical observations are frequently presented with more authority than their acquisition scope, semantics, or legal relevance can support. Privacy Lens therefore treats provenance, temporal scope, missing context, and uncertainty as first-class data. Its compliance engine rejects malformed or unsafe records, isolates overlapping evidence windows, preserves synthetic-source labels, and records the exact regulation pack and rule version applied to every finding. The evaluated legal objective is the European Union General Data Protection Regulation, but jurisdiction-specific mappings are loaded through a replaceable pack contract rather than embedded in the Android, storage, or interface layers.

The implemented product uses a three-destination information architecture comprising Overview, Findings, and Settings. It stores findings on the device, requests no sensitive runtime permission, disables Android backup, and provides explicit controls for clearing local evidence. Evaluation separates deterministic rule conformance, adversarial boundary behaviour, package integration, manifest facts, and physical-device rendering. A fixed-seed 1,000-case campaign produced 800 true positives, 200 true negatives, and no observed disagreement with the specified oracle; extended tests covered malformed inputs, temporal overlap, replay, source isolation, and regulation-pack identity. Release APK and Android App Bundle builds targeting SDK 36 succeeded, and the release APK installed and launched on an OPPO PERM00 device. These results establish a bounded engineering candidate, not legal compliance, unrestricted Android visibility, population accuracy, accessibility conformance, or application-store approval. The principal contribution is a claim-to-evidence discipline that is maintained consistently across architecture, interface, evaluation, and deployment documentation.
\end{abstract}
\newpage
\tableofcontents
\newpage
\input{Iteration-V2/Chapter1/chapter-01-10}
\section{Literature Review and Foundations}

\subsection{Regulatory Principles and Technical Signals}

The GDPR regulates processing rather than permission counts. Article 5 establishes principles including lawfulness, transparency, purpose limitation, data minimisation, accuracy, storage limitation, integrity, confidentiality, and accountability. Articles 6 and 9 concern lawful bases and special-category data; Article 25 requires data protection by design and by default \cite{eurlex2016}. EDPB guidance describes Article 25 as a continuing obligation implemented through appropriate technical and organisational measures \cite{edpb2019}. A technical observation may support an accountability inquiry, but cannot by itself demonstrate whether these legal conditions are satisfied.

This distinction determines the framework's vocabulary. A threshold exceedance is a \emph{review signal}. Missing purpose, lawful-basis, necessity, data-subject, controller, retention, and sharing evidence is shown explicitly. The UI never derives a finding labelled ``GDPR violation'' or ``compliant''. Regulation-pack references provide a reason to investigate and a link to an official source; they do not replace legal interpretation.

\subsubsection{From Legal Principle to Review Question}

Operationalisation is necessary in any regulatory technology, but it must preserve the semantic distance between a principle and a sensor. Article 5(1)(a) concerns lawfulness, fairness, and transparency. A permission count can contribute to transparency by revealing a pattern that was previously invisible, yet it cannot establish lawfulness or fairness. Article 5(1)(b) concerns purpose limitation. A count contains no purpose unless purpose is supplied as contextual evidence. Article 5(1)(c) concerns data minimisation. Frequency can be relevant to necessity, but a high frequency may be required for continuous navigation while a single collection may be excessive for an unrelated purpose. Article 5(2) concerns accountability. An audit record can support demonstration, but it is only one component of the records, decisions, safeguards, and governance required from a controller \cite{eurlex2016}.

Articles 6 and 9 further show why a permission signal cannot become a verdict. Article 6 provides several possible lawful bases; consent is not automatically the only basis. Article 9 establishes a prohibition and exceptions for special categories, but not every Android dangerous permission maps directly to special-category processing. Location, contacts, microphone, and camera data are context-dependent. Audio might contain health information, ordinary conversation, environmental noise, or no retained personal data. A technically observed invocation does not identify which of those states occurred.

The pack therefore maps a signal to a sequence of questions. For a location-related pattern, the reviewer may ask what purpose was declared, whether continuous precision was necessary, whether a less intrusive source was available, whether the application was in the foreground, whether location was retained or transmitted, and whether the user could reasonably anticipate the flow. For a microphone-related pattern, the reviewer may ask whether recording was user-initiated, how activation was signalled, whether raw audio left the device, and how long any derivative was retained. These questions are not answers; their value is that they prevent the visible count from crowding out legally decisive facts.

\subsubsection{Data Protection by Design and Interface Conduct}

Article 25 directs controllers to implement appropriate measures and privacy-protective defaults, taking account of state of the art, cost, context, and risk. EDPB guidance emphasises that this is a continuing obligation rather than a certification achieved once \cite{edpb2019}. For Privacy Lens, data protection by design affects both internal processing and the claims made to users. Local storage, disabled backup, no analytics, no advertising, no account, explicit deletion, and minimised permissions reduce the application's own data footprint. Fail-closed validation and immutable pack metadata reduce the risk that incomplete evidence becomes a misleading result.

Interface conduct is also relevant. A privacy tool could employ alarming colours, default users into sharing diagnostic data, hide evidence limitations, or imply that purchasing a service resolves risk. EDPB guidance on deceptive design patterns provides a broader reminder that formal information can still be presented in a manipulative way \cite{edpb2022}. Privacy Lens uses restrained severity language, places limitations near the result, and keeps the demonstration optional. The implementation is not claimed to have been audited against every deceptive-design criterion, but the guidance informs the design objective of calibrated rather than maximised reliance.

\subsection{Android Visibility and Permission Semantics}

Android permissions are part of a layered security model. Manifest declarations, runtime grants, AppOps state, operating-system services, and OEM behaviour do not expose one universal event stream to an ordinary third-party application. Android's permission guidance emphasises minimising permission requests and using platform-supported workflows \cite{androidpermissions2026}. Therefore, a credible review tool must state its acquisition authority and must not equate an empty result with absence of processing.

Static manifest analysis and dynamic monitoring answer different questions. Static tools can identify declared capabilities and embedded components; dynamic or taint-based systems can observe selected flows but may require instrumentation, system modification, a VPN, or elevated authority. PiOS illustrates the value and complexity of static data-flow analysis \cite{49}. Privacy Lens deliberately occupies a lighter position: it evaluates bounded audit summaries and communicates their coverage rather than claiming universal enforcement or surveillance.

\subsubsection{A Taxonomy of Mobile Privacy Evidence}

Existing approaches can be organised by the claim their evidence supports. Manifest analysis answers what capabilities an application declares and which components are packaged. It is reproducible and can be performed without executing the application, but it cannot establish that a declared capability was exercised. Static data-flow analysis attempts to connect sources and sinks in program code. It can reveal possible flows beyond manifest declarations, although reflection, native code, dynamic loading, and environment-specific paths complicate coverage. PiOS and related systems demonstrate the analytical value of tracing potential information flows while also illustrating that such analysis is specialised rather than a routine end-user feature \cite{49}.

Runtime instrumentation answers which monitored operations occurred during a particular execution. Its result depends on exercised paths, instrumentation placement, operating-system authority, and the definition of an event. Network interception answers which traffic crossed an observed interface and may expose destinations or payload structure. It cannot see purely local processing and may be limited by encryption, certificate pinning, protocol choice, and traffic that bypasses the observation path. Platform privacy dashboards can offer authoritative information for selected permissions, but their presentation and export capabilities are determined by the operating system.

Enforcement systems answer yet another question: whether a protected operation was blocked or modified under a policy. Enforcement is stronger than observation for the operation under control, but it can require device-owner status, system integration, rooting, a custom operating system, or application rewriting. These deployment choices change the threat model. A VPN monitor, for example, gains visibility by becoming an intermediary and must itself be trusted. A rooted monitor can inspect more state but weakens assumptions that apply to ordinary consumer devices.

Privacy Lens does not attempt to dominate this taxonomy. It accepts a structured summary from a bounded source and focuses on interpretation, traceability, and communication. This makes it complementary to manifest scanners, platform dashboards, enterprise management systems, and forensic monitors. A future authorised deployment could supply richer evidence through the same decision contract, but the application should continue to display the authority and coverage of that source.

\subsubsection{Historical Operations and the Ordinary-App Boundary}

The existence of an Android API is not evidence that any application can use it for any package. Access control can depend on signature permissions, usage-access grants, AppOps policy, device-owner status, process identity, Android release, and OEM modification. Historical-operation methods may be public in an SDK while useful cross-package results remain restricted in ordinary deployment. The physical-device evidence collected for this thesis exposed application-owned modes and scheduling state; it did not demonstrate unrestricted enumeration of other packages.

This boundary affects evaluation design. A complete system has an acquisition layer and a decision layer. The acquisition layer determines what the operating system or an authorised integration provides. The decision layer validates, aggregates, maps, stores, and presents that evidence. Synthetic injection can test the second layer exhaustively without proving the first. Conversely, an acquired record can demonstrate that a bridge works without proving the correctness of every decision boundary. The thesis reports these layers separately.

\subsection{Human Factors and Calibrated Trust}

Privacy notices impose substantial reading and comprehension costs \cite{12,13}. Adding another dense dashboard can reproduce the problem it intends to solve. Trustworthy warning design should therefore support quick orientation while preserving access to rationale and provenance. ISO 9241-11 frames usability through effectiveness, efficiency, and satisfaction in a specified context of use \cite{iso924111}. WCAG 2.2 and Android accessibility guidance further motivate non-colour status cues, meaningful labels, and sufficiently large touch targets \cite{wcag22,androidaccess}.

These sources support four design choices. First, the overview leads with the current evidence boundary. Second, colour is paired with iconography and text. Third, details are progressively disclosed through finding cards. Fourth, destructive actions such as clearing local findings are separated from routine navigation. These are conformance-oriented design decisions, not proof of usability; participant studies remain necessary.

\subsubsection{Cognitive Load and Progressive Disclosure}

Cognitive Load Theory distinguishes load inherent to a task from load introduced by its presentation \cite{27}. Privacy review has substantial intrinsic complexity because it combines technical events, purposes, law, expectations, and uncertainty. A product cannot remove that complexity, but it can avoid adding unnecessary navigation, decoration, or jargon. The earlier prototype mixed health dashboards, research projects, experiment controls, scoring views, and privacy findings. That structure made the user's immediate task unclear. The three-destination Revision 10 architecture reduces extraneous choice and uses a stable pattern for each finding.

Progressive disclosure is not equivalent to hiding caveats. The most important caveat---the limited reach of evidence---appears on the overview before the user initiates a review. Detail cards then reveal the observed count, source, pack, reference, rationale, and missing context. This ordering allows rapid orientation while preserving auditability. A legal citation without an explanation would be too terse; a full reproduction of statutory text on every card would be too heavy. The interface therefore supplies a short reference and offers the official source separately.

Dual-process accounts distinguish fast, heuristic judgement from slower deliberation \cite{22}. Severity colour can produce a rapid response, which is useful for prioritisation but dangerous if colour appears to communicate legal certainty. Textual labels, source chips, explicit caveats, and neutral action language are intended to interrupt an overly simple red-equals-illegal inference. This intention requires empirical testing. The thesis does not infer comprehension from the presence of these elements.

\subsubsection{Trust Calibration and Explanation Quality}

Trust is calibrated when reliance corresponds to actual capability. A system that understates uncertainty can induce over-reliance; a system that only displays caveats may become unusable and be ignored. Lee and See describe the importance of appropriate reliance on automation \cite{34}. In Privacy Lens, calibrated trust is pursued through claim symmetry: a positive-looking result and a warning must both show their evidence basis. A no-technical-concern state means that no rule fired within the admitted evidence, not that processing was lawful. A needs-review state means that a rule fired, not that a violation occurred.

Explanation quality has at least four dimensions. Fidelity asks whether the explanation corresponds to the implemented rule. Completeness asks whether the important inputs and limits are represented. Comprehensibility asks whether the intended audience can understand it. Actionability asks whether the explanation supports a reasonable next step. The deterministic rule and stored pack metadata support fidelity. Missing-context questions support completeness. The card hierarchy aims at comprehensibility. Suggested review questions support actionability. Only the first two dimensions are strongly evaluated in the present work; comprehension and actionability need participant evidence.

\subsubsection{Accessibility as an Evidence Requirement}

Accessibility affects whether transparency is available in practice. WCAG 2.2 requires that information is not conveyed through colour alone and includes criteria concerning contrast and target size \cite{wcag22}. Android guidance recommends meaningful content descriptions and interaction targets that are large enough to operate reliably \cite{androidaccess}. The implementation pairs icons and words with colour and supplies accessibility labels for primary controls.

These measures are design evidence, not conformance certification. True accessibility evaluation should include screen-reader order, announcements after asynchronous review, large-font reflow, switch access, keyboard focus, high-contrast settings, reduced motion, colour-vision deficiencies, and comprehension with users who rely on assistive technology. The distinction mirrors the wider thesis: source inspection can establish that a label exists, whereas a user study is needed to establish that the resulting interaction is effective.

\subsection{Replaceable Rule Systems}

Hard-coding statutory names into a decision engine creates maintenance and validity risks. A legal text may change, another jurisdiction may use different concepts, and a translated label may not preserve the same meaning. The appropriate software pattern is a stable engine contract with versioned, independently reviewable rule data. A pack should declare its jurisdiction, version, official source, caveat, supported signals, thresholds, references, and explanatory text. The engine should reject unknown packs and preserve the selected pack identifier in every finding.

Replaceability does not mean equivalence. Two packs cannot be treated as interchangeable simply because they implement the same TypeScript interface. Each requires independent legal review, localisation, validation fixtures, and governance over updates. The included ``Research Baseline'' pack demonstrates switching without representing a second jurisdictional legal interpretation.

\subsubsection{Pack Lifecycle and Review Governance}

A credible pack lifecycle begins before implementation. The maintainer identifies the intended jurisdiction, material and territorial scope, regulated actors, effective date, official consolidated source, language versions, and the technical signals actually available. A qualified reviewer then maps signals to questions, not verdicts, and records assumptions and exclusions. Engineering converts the reviewed mapping into immutable data and fixtures. Localisation is reviewed separately because legally important distinctions can be lost even when the code remains valid.

Release governance should treat a pack as a versioned dependency. A substantive legal amendment, authoritative interpretation, corrected translation, changed threshold, or altered missing-context question may require a new pack version. Findings retain the version used at evaluation time. Updating the active pack should not rewrite historical records. If a migration is needed, it should create a new evaluation linked to the old one rather than silently changing the original rationale.

Pack selection also raises a scope question. Device location alone is not a reliable selector because applicable law can depend on establishment, targeting, residence, processing context, contracts, and sector. Revision 10 therefore requires explicit user selection from registered packs and presents the jurisdiction label and caveat. Automatic selection could be added only with a defensible scope model and an override.

\subsubsection{Rule-Based and Learned Decision Models}

Rule-based systems are attractive for a legal-support interface because a reviewer can reconstruct the condition that fired. A fixed threshold also has obvious limitations: it can be bypassed by splitting activity across windows, may not represent purpose, and may treat different application categories alike. Learned anomaly detection can adapt to patterns, but it introduces training-data bias, distribution shift, opacity, and uncertainty. A hybrid system could combine transparent guardrails with contextual or personalised models, provided that model provenance and uncertainty remain visible.

The present engine intentionally uses explicit research thresholds. The contribution lies in safe admission, temporal isolation, versioning, explanation, and boundary disclosure rather than in claiming optimal thresholds. Population calibration would require representative applications, ground-truth definitions, stratification by purpose and category, independent annotation, and validation on data not used to choose parameters. Without that evidence, a mathematically sophisticated adaptive threshold could be less trustworthy than a simple rule whose limits are stated.

\subsubsection{Automated Validation and Oracle Independence}

Randomised tests can explore more input combinations than a short hand-written suite, but their evidential value depends on the generator and oracle. Uniform random counts may seldom exercise exact thresholds. A generator that only creates well-typed values cannot test the runtime boundary. An oracle copied from the production helper may reproduce the same defect. The evaluation therefore combines fixed seeds, boundary values, malformed types, permission mutation, temporal constructions, and an explicit expected-result calculation.

Fuzzing and property-oriented testing are especially useful for invariants. Counts should remain finite and non-negative. Results should not depend on the order of unrelated package histories. Replaying an identical window should not increase totals. Overlapping aggregate windows should not be added. A simulator source should never become a native-observation source. Unknown packs should not load. These properties express safety behaviour more clearly than a single accuracy number.

Statistical language must also match the sample. A deterministic 1,000-case run provides strong regression coverage for its generator, but it is not a random sample of real applications or alleged infringements. Confidence intervals applied to such a run would describe the generated distribution, not the mobile ecosystem. Revision 10 reports the confusion matrix descriptively and reserves population claims for future independent datasets.

\subsection{Store Delivery as a Trust Boundary}

Application-store readiness includes more than a successful APK. Google Play uses Android App Bundles for new application delivery, requires current target API levels, and requires developers to complete a Data safety declaration \cite{androidbundle2026,playtarget2026,playdatasafety2026}. Signing identity, a public privacy-policy URL, support contact, screenshots, listing text, and console review remain owner-controlled steps. A research build can be an initial submission candidate while still being unpublishable until these external obligations are completed.

\subsection{Synthesis}

The literature yields four gaps addressed by the prototype: permission signals are easily over-interpreted; acquisition coverage is often hidden; legal mappings are commonly coupled to product code; and engineering success is frequently conflated with deployment or legal validity. The following chapter converts these gaps into explicit system requirements and evidence contracts.

\input{Iteration-V2/Chapter3/chapter-03-10}
\section{Implementation}

\subsection{Application Identity and Delivery Configuration}

Revision 9 presents the prototype as \emph{Privacy Lens}, version 1.1.0 (version code 2), under the Android application identifier \texttt{com.zhihengzhang.privacylens}. The Android configuration uses minimum SDK 24 and target SDK 36. Release builds generate both an APK for controlled installation and an Android App Bundle for store delivery. The Gradle signing configuration accepts upload-key values through environment variables; when they are absent, local quality-assurance builds use the debug signing fallback and are not represented as production-signed artefacts.

Android backup is disabled. The manifest explicitly removes legacy external-storage permissions. The merged release manifest contains Internet access for official-source links and the operational permissions introduced by WorkManager for wake locks, network state, boot rescheduling, and foreground-service compatibility. It requests no sensitive runtime permission.

\subsection{Rule-Pack Engine}

The generic rule-pack engine implements validation and threshold evaluation against a supplied regulation definition. The earlier GDPR-named engine remains as a compatibility wrapper that delegates to the EU pack; it no longer owns the rule table. The registry exports only recognised packs and rejects an unregistered identifier.

The EU pack maps supported permission signals to review rationales and official GDPR references. Its caveat states that output is a technical review aid. The Research Baseline pack uses deliberately non-legal terminology. Both conform to the same interface, allowing the Settings screen to change the active pack without changing bridge, repository, dashboard, or navigation code. Pack identity is persisted with every finding, so switching the active selection does not relabel historical evidence.

\subsection{Validation and Temporal Isolation}

The engine validates package name, permission key, source, count, window, timestamp, and context before evaluation. JavaScript unsafe integers, negative counts, reversed windows, unknown sources, and malformed identifiers fail closed. Simulator provenance is an enumerated source rather than a boolean UI decoration.

The repository uses a versioned local schema and migrates earlier records. Exact replay windows replace prior values. Overlapping aggregate windows cannot be summed; adjacent completed windows remain eligible for a bounded rolling signal. Records are partitioned by package and permission and capped per partition. These constraints make the decision reconstructable while acknowledging that aggregate summaries lack event identifiers.

\subsection{Native Android Boundary}

Kotlin components implement the React Native bridge, scheduling, permission-audit worker, and controlled violation simulator. The bridge exposes capability status so that the interface can report whether native audit support is present. The product does not claim unrestricted access to other applications' AppOps history. A blank audit is displayed alongside an evidence-reach warning, because stock Android visibility depends on deployment authority and OEM behaviour.

The simulator produces synthetic records for demonstration and regression. Its data source remains explicitly labelled through storage and presentation. This enables a user to inspect a populated workflow without confusing the output with observations collected from the device.

\subsection{Information Architecture and Visual Design}

The interface was reduced to three primary destinations:

\begin{itemize}
    \item \textbf{Overview} explains the product purpose, initiates an audit, summarises finding classes, and states the evidence-reach boundary;
    \item \textbf{Findings} separates review signals, evidence gaps, and no-concern results, with source and pack labels on each card;
    \item \textbf{Settings} selects a registered rule pack, explains its caveat and official source, runs the synthetic demonstration, and clears local findings.
\end{itemize}

A custom bottom navigation component replaced the earlier crowded tab structure. The visual system uses deep evergreen, teal, mint, off-white, restrained risk accents, rounded cards, and generous whitespace. The icon combines a shield and lens aperture without text. Status is carried by words and icons as well as colour. Buttons and navigation targets use large touch areas and meaningful accessibility labels, following Android and WCAG design guidance \cite{androidaccess,wcag22}.

\subsection{Privacy and User-Control Implementation}

Findings and active-pack choice use app-private local storage. There is no account, analytics SDK, advertising SDK, remote evidence database, or upload workflow. Settings provides a direct clear-local-findings action. The repository includes a bilingual privacy-policy draft and a store-readiness checklist. The policy distinguishes application behaviour from owner actions still required before public release, including a stable HTTPS policy URL and support contact.

\subsection{Codebase Reconciliation}

Legacy health-dashboard, project, scoring, experiment, charting, duplicate compliance-service, and unused template components were removed. Dependencies for removed features were also deleted. This cleanup reduces navigation ambiguity and prevents unsupported research screens from appearing to be finished product functions. The resulting codebase has one compliance engine path, one privacy context, one navigation model, and one current user guide.

\subsection{Technology Stack and Module Boundaries}

The application uses Expo and React Native for the interface, TypeScript for the decision and state layers, AsyncStorage for app-private persistence, and Kotlin for the Android bridge and workers. Expo Router supplies route composition, while React Navigation primitives support the custom bottom navigation. AndroidX WorkManager represents deferrable background work. The release build uses the Android Gradle Plugin, Hermes JavaScript engine, and standard Android packaging tasks.

The stack was selected to separate concerns rather than to maximise the number of frameworks. TypeScript provides shared types for evidence, findings, and packs. Kotlin is limited to capabilities that require Android APIs or native scheduling. The interface does not import Kotlin details; it calls a typed service that reports capability and results. The compliance engine does not import React Native, storage, or Android modules. This makes deterministic tests runnable under Node after TypeScript compilation.

At the source-tree level, regulation definitions reside under the regulation directory, generic decision logic under compliance, application state under context, the bridge adapter under services, and reusable presentation components under a privacy-specific component directory. Theme constants are centralised. This organisation replaces earlier duplicate compliance services and prevents a screen from becoming the authoritative home of a rule.

\subsection{Regulation-Pack Implementation}

The EU pack is an immutable TypeScript object identified as \texttt{EU\_GDPR}. It declares the European Union and EEA as its jurisdiction label, cites Regulation (EU) 2016/679 as its version label, and links to the consolidated EUR-Lex source. The pack includes a legal caveat that the automated warning supports accountability review and does not determine infringement.

Three permission families are implemented. Location has a research baseline of 24, microphone a baseline of 8, and contacts a baseline of 4. Each baseline uses a multiplier of 1.5, producing prompt thresholds of 36, 12, and 6 respectively. These values are prototype parameters. They are not extracted from GDPR text, not population estimates, and not represented as legal limits. Their purpose is to exercise an explainable evaluation path and provide deterministic boundaries for testing.

The EU pack supplies a rationale and reference for every permission family. Location directs review toward necessity, proportionality, and documented lawful basis under Articles 5 and 6. Microphone review includes the possibility that special-category material may require Article 9 analysis. Contacts review emphasises purpose limitation, minimisation, and lawful basis. The reference is intentionally a prompt rather than an assertion that the cited provision applies to every record.

Missing-evidence logic is also pack-owned. The EU pack asks for specified purpose, Article 6 lawful basis, controller identity, retention period, and, when special-category processing is marked, an Article 9 condition. If processing relies on consent, consent has been withdrawn, and a positive access count remains, the pack returns its most urgent review status. If required context is missing, it returns insufficient evidence. Otherwise a threshold signal produces review required, while the absence of a signal produces no technical concern for the admitted check.

The Research Baseline pack uses higher prototype baselines of 36, 12, and 6 before applying the same multiplier. It uses region-neutral principles and does not require an Article 6 or Article 9 field. Its caveat explicitly states that it is not law or legal advice. The difference demonstrates that the engine can consume pack-owned rules and missing-evidence logic without embedding GDPR branches in generic code.

\subsection{Runtime Validation Pipeline}

The safe evaluation entry point accepts an unknown value rather than trusting its compile-time annotation. The first check requires a non-array object. The package identifier must be a string between one and 255 characters using a conservative character set. Permission type must belong to the owned set of location, microphone, and contacts. Source, if present, must be simulator, native bridge, or imported. This sequence prevents later operations from assuming properties that have not been established.

The access count must be a non-negative safe integer. The window start and end must also be safe integers, the start must precede the end, duration cannot exceed one day, and the end cannot be implausibly far in the future. When event timestamps are supplied, their array length must equal the access count and every timestamp must lie within the audit window. These conditions ensure that peak-rate calculations do not operate on inconsistent evidence.

Processing context is validated field by field. Purpose, controller identity, and Article 9 condition are optional strings. Consent-withdrawn, special-category, and user-initiated values are optional Booleans. Lawful basis must belong to the Article 6 enumeration used by the prototype. Retention days must be a non-negative safe integer. Values are rejected rather than coerced because coercion can change meaning at a trust boundary.

Every rejection returns an explicit code and message with a rejection timestamp. The throwing evaluation method is retained for callers that require an exception, but the application can use the safe result to display a controlled evidence gap. This dual interface allows strict internal use without forcing the UI to recover from an untyped exception.

\subsection{Signal Computation}

After admission, the engine loads the rule for the active pack and computes the prompt threshold

\begin{equation}
T_p=\max\left(1,\left\lceil B_pM_p\right\rceil\right),
\end{equation}

where \(B_p\) is the pack baseline and \(M_p\) is the deviation multiplier. Daily excess is \(\max(0,x-T_p)\), where \(x\) is the admitted count. A daily-total signal is emitted only when excess is positive, so equality with the threshold remains below the signal boundary.

When event timestamps are present, the engine sorts them and evaluates a sliding sixty-second interval. Two indices move monotonically through the sorted array, so peak calculation is linear after sorting. Permission-specific burst limits are 12 for location, 6 for microphone, and 4 for contacts. A burst signal is emitted when the observed peak is strictly greater than the configured limit.

The history key combines package and permission. Previous entries outside the one-day watermark are discarded. An entry with the same start and end is treated as replay and replaced. Only entries ending no later than the current start contribute to the completed rolling total. Overlapping entries stay in bounded history but do not contribute to that total. A cross-window signal is disabled for simulator evidence and requires at least one compatible completed window.

Risk derives from the greatest normalised excess across daily, burst, and rolling dimensions. Ratios below the first threshold remain low; larger ratios map through medium, high, and critical bands. Risk controls presentation and notification priority, while the pack separately controls the compliance-review status. Keeping these concepts distinct prevents a large technical deviation from bypassing missing legal context.

\subsection{Finding Construction and Communication}

The finding identifier combines package and permission so the repository can update a stable subject. Evidence fields contain daily count, peak calls per minute, rolling count, and source. Decision fields contain threshold, excess count, active signals, risk, detection time, regulation identifier and name, legal or research reference, and rationale. Compliance fields contain status, applicable principles, missing evidence, and caveat.

Communication is derived after the finding exists. The title identifies the permission review. The summary translates status and signal identifiers into human-readable language. The recommended action first requests missing evidence; an urgent result suggests pausing processing where appropriate and obtaining human review; other results recommend examining purpose, necessity, and proportionality. Notification priority is urgent only for the strongest status or critical risk, silent for no technical concern, and standard otherwise.

This construction ensures that UI wording is based on stored decision state. A card is not permitted to infer a legal label merely from colour or excess count. It receives the pack name, source, caveat, and missing context that were produced with the finding.

\subsection{Repository Behaviour}

The repository uses a versioned AsyncStorage key. Listing parses the stored array, replacement writes at most the bounded presentation set, clearing removes the key, and saving updates an existing package-permission finding or appends a new one. The latest application context also normalises records created by earlier schemas so that pack identity and evidence fields remain available where they can be reconstructed safely.

Notification state is based on change. A new finding can notify; a reversal of active state can notify; an increase in risk rank can notify. A byte-for-byte identical record is not treated as changed. This policy prepares the product for attention-aware notification without claiming that a longitudinal notification campaign has been conducted.

The decision engine's internal temporal history and the persistent finding repository serve different purposes. Engine history supports signal computation within a running context. Repository findings support user review across sessions. Conflating them would either store unnecessary event aggregates indefinitely or make visible findings disappear whenever the process restarts.

\subsection{Android Bridge and Scheduling}

The Kotlin module reports whether the native privacy inspector is available and exposes bounded audit and demonstration operations to React Native. WorkManager workers are registered for periodic or one-time tasks according to Android's scheduling model. WorkManager is deferrable: the operating system may delay work according to battery, idle state, quotas, and vendor policy. The implementation therefore does not describe a registered interval as proof of exact execution time.

The audit worker reads only evidence available under the application's deployment authority. On the tested ColorOS device, that authority did not expose unrestricted third-party history. The bridge returns capability and available summaries rather than fabricating other-package events. The simulator worker follows a controlled path and marks every generated record as simulator source.

Native package declarations were moved from the anonymous template namespace to the stable Privacy Lens namespace. Activity, application, bridge, package, and worker source files use the same identifier as Gradle and the manifest. This reconciliation matters for release upgrades, signing continuity, deep links, provider authorities, and store identity.

\subsection{Interface Implementation}

The Overview screen has four layers. A brand header establishes identity and provides a direct Settings entry. A high-contrast hero panel explains the evidence-before-conclusions proposition and presents the primary review action. A current-review section summarises needs-review, evidence-gap, and no-concern counts. An evidence-reach card explains whether the native bridge is available and why an empty result is not proof of no access.

The Findings screen is designed for both empty and populated states. The empty state explains what will appear and directs the user back to a review or demonstration. The populated state renders reusable finding cards. Each card shows status, permission, package, source, regulation pack, detected signals, rationale, and missing context. Source and regulation are placed near the result because they alter how the result should be interpreted.

The Settings screen separates preference, explanation, demonstration, and deletion. Pack selection is restricted to registry entries. Selecting a pack displays its jurisdiction, version, description, caveat, and official source. The controlled demonstration is visibly separated from device review. Clearing findings uses a destructive visual treatment and explanatory copy.

The custom bottom navigation uses three equally weighted targets with icon, label, accessibility role, selected state, and sufficient touch area. It replaced the platform tab implementation after physical coordinate and hierarchy evidence revealed ambiguity in the previous navigation setup. Navigation semantics no longer depend on obsolete routes.

\subsection{Brand and Accessibility Implementation}

The visual language uses evergreen for trust and stability, teal for actions, mint for supportive surfaces, warm off-white for the application background, and restrained red and blue accents for review and evidence states. Rounded panels group related information without relying on heavy separators. Typography uses strong size contrast between product proposition, section heading, count, label, and explanatory copy.

The application icon combines a shield with a lens aperture. The shield suggests protection while the aperture suggests inspection; neither implies surveillance of a person. Adaptive foreground and monochrome assets support launcher contexts. Splash imagery, application icon, colours, and product name are generated consistently across Android densities.

Accessibility labels are supplied for navigation and important actions. Icons do not carry status alone. Buttons use text and large target areas. Cards maintain text contrast against their surfaces. These implementation facts support an accessibility-oriented design claim, while screen-reader traversal, large-font reflow, and assistive-technology user testing remain future acceptance work.

\subsection{Build, Permission, and Signing Configuration}

The Android build sets minimum SDK 24, compile and target SDK 36, version code 2, and version name 1.1.0. Release tasks produce an APK for controlled device installation and an AAB for Play delivery. The build reads upload-store path, passwords, and alias from environment variables. If they are absent, a local QA fallback signs with the debug configuration. The build can therefore be tested without placing secrets in source, while its output remains clearly distinguished from a production upload artefact.

The source manifest requests Internet access and explicitly removes inherited legacy external-storage permissions using manifest-merger directives. The final merged manifest also contains operational declarations contributed by WorkManager and AndroidX. Android backup and data extraction are disabled for the application. Inspection of the installed release package confirms that no sensitive runtime permission is requested.

Package reduction was part of implementation. Dependencies used only by removed charts, health views, haptics, image helpers, secure storage, web browser wrappers, gestures, SVG rendering, and network-status features were removed when no longer needed. Generated build and cache directories remain ignored. This reduces attack surface and maintenance burden, although transitive dependencies are still reviewed through the lock file and final manifest.

\subsection{Implementation Summary}

The implemented system realises the Chapter 3 contracts: inputs fail closed, evidence provenance survives the pipeline, packs are replaceable through a registry, the interface avoids legal verdict language, and Android delivery artefacts use a stable identity. Chapter 5 evaluates these properties at source, rule, build, manifest, and physical-device layers.

\section{Evaluation and Results}

\subsection{Evidence Model}

The evaluation separates five evidence layers: source conformance, deterministic rule behaviour, adversarial boundary behaviour, Android package integration, and physical-device usability evidence. This separation prevents a passed unit test from being reported as device reliability or legal validity. The final acceptance record is stored at \path{testing-report/release-candidate-2026-08-09.md}; screenshots and UI hierarchies are retained in the adjacent dated folder.

\subsection{Source and Rule Verification}

The complete application passed \texttt{tsc --noEmit} and ESLint. The compliance sources were compiled before execution, preventing stale generated JavaScript from being treated as current evidence.

A fixed-seed logical campaign generated 1,000 cases. Eight hundred were expected review signals and 200 were below-threshold controls. The observed confusion matrix was

\begin{equation}
(TP,TN,FP,FN)=(800,200,0,0),
\end{equation}

which gives precision and recall of 1.0000 against the specified oracle. The oracle and engine implement the same published threshold model; this is a conformance and regression result, not an independent estimate of real-world infringement detection. Zero observed errors do not establish a zero population error rate.

Extended suites passed for unsafe types and integers, malformed context, unknown source, permission-key mutation, temporal overlap, adjacency, replay, retention, simulator isolation, dashboard presentation, missing-context disclosure, and regulation-pack identity. The tests confirm the enumerated boundaries. They cannot establish correctness for telemetry the application never receives or for legal interpretation outside the fixtures.

\subsection{Release Build and Manifest Verification}

Gradle completed \texttt{assembleRelease} and \texttt{bundleRelease} with target SDK 36. Table~\ref{tab:r9-artifacts} records the generated artefacts.

\begin{table}[H]
\centering
\small
\caption{Revision 9 Android release artefacts}
\label{tab:r9-artifacts}
\begin{tabular}{p{0.31\textwidth}r p{0.42\textwidth}}
\toprule
Artefact & Bytes & SHA-256 \\
\midrule
\texttt{app-release.apk} & 62,762,819 & \texttt{\seqsplit{4A5BCED7FBF6A7A5A0E337BFAADB003A3C7FE89FDAFBF9BADCBE6BC4CBFF51B5}} \\
\texttt{app-release.aab} & 31,700,025 & \texttt{\seqsplit{E4090FEEF4EFBD0B697F863D80D28F289C3554BA89394A1435EACA2C226ACF78}} \\
\bottomrule
\end{tabular}
\end{table}

The merged release manifest contained declarations for Internet access, wake lock, network state, boot rescheduling, foreground service, and an app-scoped signature permission. External-storage read and write permissions were absent. Inspection of the installed package confirmed the same requested-permission set and no runtime permissions. Android backup was disabled in the application declaration.

\subsection{Physical-Device Evaluation}

The release APK was installed by ADB on one OPPO PERM00 ARM64 device. The installed package reported \texttt{com.zhihengzhang.privacylens}, versionName 1.1.0, versionCode 2, minSdk 24, and targetSdk 36. Launcher cold start succeeded. A filtered scan of recent \texttt{AndroidRuntime} and \texttt{ReactNativeJS} error channels returned no fatal match during the acceptance run.

The three primary destinations rendered at the device's captured application resolution. Coordinate-driven exploration initially appeared to show a navigation-target problem because the physical display and captured application coordinate spaces differed. A bounded coordinate sweep and UI-tree inspection identified the offset; subsequent navigation and custom bottom-bar interaction succeeded. This calibration episode is retained because it demonstrates why raw ADB coordinates must not be interpreted without the target's current window geometry.

The controlled demonstration completed a 50-record synthetic workflow. The populated overview displayed four review items, zero evidence gaps, two no-concern items, and descriptive precision and recall of 100\% for the simulator's known labels. The Findings screen marked the source as synthetic and displayed the EU GDPR pack. These values validate the demonstration path only; they are not device-observation accuracy.

\subsection{Usability and Trust Assessment}

The product was assessed against three release-candidate questions.

\begin{table}[H]
\centering
\small
\caption{Bounded product assessment}
\begin{tabular}{>{\raggedright\arraybackslash}p{0.19\textwidth}>{\raggedright\arraybackslash}p{0.32\textwidth}>{\raggedright\arraybackslash}p{0.37\textwidth}}
\toprule
Criterion & Supporting evidence & Boundary \\
\midrule
Visually coherent & consistent brand, hierarchy, spacing, cards, iconography, and three-destination navigation on the tested device & one device; no participant preference study \\
Trustworthy & local-first disclosure, source and pack labels, explicit missing evidence, no legal-verdict language & not a legal audit or security certification \\
Willing to use & clear first action, concise summaries, optional details, demonstration, and direct data clearing & evaluator judgement; no longitudinal adoption measure \\
\bottomrule
\end{tabular}
\end{table}

The evaluator judged the candidate to meet the project's ``beautiful, trustworthy, and willing to use'' gate for an initial submission candidate. This is a structured expert judgement supported by rendered evidence, not a statistically generalisable user study.

\subsection{Evaluation Questions and Acceptance Logic}

The evaluation was organised around six questions. First, does the source compile and satisfy static rules? Second, does the decision engine agree with its specified thresholds and statuses? Third, does it reject malformed or adversarial evidence safely? Fourth, do regulation packs remain independent and traceable? Fifth, can the Android release artefacts be assembled, inspected, installed, and launched? Sixth, does the rendered product communicate the expected hierarchy and evidence boundaries on the available physical device?

Each question has a distinct acceptance condition. Source acceptance requires a clean TypeScript application check and ESLint run. Rule acceptance requires exact expected confusion-matrix values for the fixed corpus and passing boundary assertions. Runtime-boundary acceptance requires every malformed case to return the specified rejection code without an uncaught exception. Pack acceptance requires registry allowlisting, different thresholds where defined, retained pack identity, and a non-legal caveat for the Research Baseline.

Android build acceptance requires both release tasks to complete and the artefacts to exist with recorded hashes. Permission acceptance requires agreement between the merged manifest and installed package, with legacy external-storage permissions absent. Physical acceptance requires successful streamed installation, launcher start, stable primary navigation, captured final states, and no matching fatal error in the bounded log scan. Visual acceptance requires inspection of actual screenshots rather than source code alone.

The protocol deliberately does not set an acceptance condition for GDPR compliance, legal accuracy, population detection accuracy, long-run scheduler reliability, battery consumption, multi-device compatibility, or store approval. Those outcomes require evidence not collected in this campaign.

\subsection{Reproducibility Controls}

The compliance sources are compiled immediately before the generated JavaScript runners execute. This control was introduced because a failed TypeScript command can otherwise leave older generated output in place, producing a misleading green test. The acceptance sequence checks the compiler exit code before invoking Node. Generated compliance output is ignored by Git so it cannot become an accidental source of truth.

The main logical corpus uses a fixed linear-congruential seed. A fixed seed makes a failure reproducible while still exploring varied generated values. The 1,000 rounds preserve an 80/20 design: four of every five cases are expected review signals and one is a control. A fresh engine is created for each main-corpus case so temporal state from one generated subject cannot leak into another.

Boundary cases are not left to random chance. For every supported permission, counts at threshold minus two, minus one, equal, plus one, and plus two are exercised. The assertion encodes the strict boundary: only values above the prompt threshold become active under the daily-total rule. Burst tests similarly check exactly at and exactly above each permission-specific limit.

Time is fixed in extended tests. Windows and timestamp arrays are created relative to a constant epoch value rather than the wall clock. This prevents a test from becoming invalid while it runs and makes future failures comparable. Production validation still compares implausible future times against the current clock, but fixtures choose values far enough in the past to avoid accidental rejection.

\subsection{Deterministic Corpus Interpretation}

The confusion matrix compares the engine's active-state decision with the generator's specified expected label:

\begin{equation}
\mathrm{Precision}=\frac{TP}{TP+FP}, \qquad
\mathrm{Recall}=\frac{TP}{TP+FN}.
\end{equation}

The resulting precision and recall of 1.0000 mean that the implementation agreed with the oracle for all generated cases. This supports a regression claim: later refactoring did not alter the expected threshold semantics for the corpus. It also supports a fixture-quality claim: the corpus contains both positive and negative cases and the metric implementation returns the expected counts.

It does not show that a threshold crossing identifies unlawful processing. The expected label is generated from the same documented research model that the engine is meant to implement. Although the oracle is expressed independently of the evaluated result, both are derived from one specification. A defect in that specification could be reproduced consistently. There is no independently labelled sample of real applications or legal cases.

The tests also verify metric edge conditions. Empty samples, all-negative samples, and samples with missed positives must produce finite precision and recall rather than not-a-number or infinity. This matters because a dashboard can otherwise display an apparently precise but meaningless metric when a denominator is zero.

\subsection{Adversarial Input Results}

The initial security cases reject unsupported permission values, not-a-number counts, negative counts, reversed windows, and prototype-like permission keys. Extended cases add unsafe integers and positive infinity. Every such count returns the invalid-count code. Timestamp evidence is rejected when array length differs from count, when an element falls before the start, when it falls after the end, or when it is fractional.

Processing context is attacked through type confusion. Numeric purpose, array controller identity, object Article 9 condition, string consent-withdrawn, numeric special-category, and string user-initiated values all return invalid context. An unknown lawful basis and negative retention period are also rejected. These cases are significant because JavaScript truthiness or an early string operation could otherwise convert malformed evidence into a plausible status or throw outside the controlled error path.

Source validation rejects a forged source. This prevents a caller from using an arbitrary label that resembles an authoritative acquisition path. Package and permission history are isolated: a high-count record for one package cannot make another package active. Histories are also isolated by permission family.

The assertions demonstrate the implemented rejection surface, not complete security. They do not cover a compromised native module, a modified application binary, corrupted AsyncStorage beyond the migration parser, malicious operating system, or every possible Unicode identifier. Security conclusions remain bounded to the enumerated input contract.

\subsection{Temporal and Replay Results}

Burst tests create timestamp arrays exactly at and just above the location, microphone, and contacts limits. At-limit arrays produce no burst signal; one additional event produces the signal. A separate imported-evidence example concentrates 20 location events within one minute and is detected.

Cross-window tests use two below-threshold location summaries. The first remains inactive. When the second begins after the first has completed, their rolling total exceeds the threshold and the engine emits a cross-window signal. The same construction is repeated with adjacent windows to establish that equality between a prior end and current start is admissible.

An overlapping construction uses one interval that intersects the next. The second finding's rolling count remains its own count and no cross-window signal appears. This is the expected conservative behaviour because aggregate evidence cannot identify duplicate events in the overlap. An exact replay is then evaluated twice before a compatible next interval. Rolling total confirms that the repeated window contributed once, not twice.

Simulator histories are excluded from cross-window signalling. The demonstration may still produce a daily or burst signal for a record, but it cannot present accumulated simulator activity as imported temporal evidence. This is a provenance safeguard, not a statistical assumption.

The temporal tests operate in process memory and do not demonstrate persistence of aggregate history across process death. They establish the logic of the current engine instance. A production design requiring durable cross-window evidence would need a versioned evidence-history repository and migration tests separate from the visible finding repository.

\subsection{Accountability and Presentation Results}

A low-count contacts record without processing context returns insufficient evidence rather than no concern. Its communication priority is standard because the absence of legal context warrants review. A low-count location record with purpose, contract basis, controller identity, one-day retention, and user-initiated context returns no technical concern and a silent priority.

A microphone record marked as consent-based, special-category, and accessed after consent withdrawal returns the urgent review status. The fixture does not prove a real GDPR infringement; it confirms that the pack's stated escalation branch is implemented. A special-category microphone record without an Article 9 condition lists that condition among missing evidence. An invented lawful basis and invalid retention value are rejected at admission.

Dashboard projection preserves three different states from the accountability fixtures. Summary totals show one action item, one evidence gap, and one no-concern result. Filters surface the expected finding. Every presentation state has both a textual label and marker, ensuring that source-level semantics do not depend on colour alone.

\subsection{Regulation-Pack Results}

The registry exposes at least two packs and recognises only the two expected identifiers. A forged identifier is rejected. The same location audit with a count of 40 is evaluated under both packs. The resulting findings retain different regulation identifiers and different thresholds because each pack supplies its own baseline.

This test is stronger than checking that a dropdown changes text. It passes the same admitted evidence through two independently instantiated engines and inspects the resulting finding objects. The Research Baseline result also contains a caveat stating that it is not law. The test does not validate either pack's legal interpretation; it validates decoupling, identity, and disclosure.

Future pack validation should add a contract suite automatically applied to every registry entry. It should require a unique identifier, version, official-source policy, complete supported-rule set, positive finite baselines and multipliers, non-empty rationale and caveat, deterministic missing-evidence logic, and fixtures reviewed by the responsible domain expert.

\subsection{Build Reproduction}

The release build initially encountered Windows path and dependency-layout problems. A hoisted local dependency layout allowed native CMake paths to remain within workable bounds. React Native's new architecture was disabled for the release candidate because enabling it in this Windows checkout reproduced path-related native build failures. This is a build-environment decision, not evidence that the legacy architecture is preferable for future releases.

Gradle reported deprecated features that will require attention before Gradle 9 compatibility. It also reported an SDK XML version warning caused by differences between installed Android Studio and command-line components. Neither warning prevented the recorded release build, but both are retained as maintenance debt rather than suppressed.

The release task generated bundles for multiple Android ABIs. The APK was installed on the ARM64 physical device. The AAB was not uploaded to Play Console, processed by Play signing, or delivered through an internal track. Its success establishes local bundle generation only.

Hashes were recorded immediately after the final permission-manifest rebuild. Any source or build-tool change requires rebuilding and computing new hashes. The files are generated artefacts and are not treated as immutable simply because their names remain the same.

\subsection{Manifest and Installed-Package Reconciliation}

Source-manifest inspection alone was insufficient because dependencies contribute declarations through manifest merging. The initial merged manifest exposed legacy external-storage permissions inherited from transitive Expo components. Revision 10 added explicit merger removals for read and write external storage, rebuilt both artefacts, and inspected the merged result again.

The final merged manifest contained Internet, wake-lock, network-state, boot-completed, foreground-service, and an application-scoped dynamic-receiver signature permission. These operational declarations are consistent with linking and WorkManager components. The installed-package dump reported the same requested set and no runtime permission grants. This two-level check supports the documentation statement that no sensitive runtime permission is requested.

The package dump also confirmed application identity, version name, version code, minimum SDK, target SDK, primary ABI, install state, and signing version. The release used the QA fallback signing identity because production environment variables were absent. A successful package inspection is therefore not proof of final upload-key configuration.

\subsection{Physical Interaction Protocol}

The device was identified by ADB before installation. The final release APK was streamed with replacement enabled, the package was force-stopped, and the launcher intent was injected. Package metadata was read after installation so that a stale or differently named package could not be mistaken for the new build.

Interaction evidence included screenshots and UI hierarchy dumps for overview, findings, and settings. A controlled demonstration populated the dashboard and finding list. The final overview was captured again after the permission-manifest rebuild and installation. Runtime review used filtered AndroidRuntime and ReactNativeJS error channels; no matching fatal output appeared in the bounded recent log window.

An important measurement issue arose from coordinate spaces. The physical display reported 1080 by 2400 pixels, while the captured application window and hierarchy used 1080 by 2244. Initial taps derived from the full display therefore appeared to miss navigation targets. A bounded coordinate sweep and hierarchy inspection established the offset, after which navigation succeeded. The event is reported because excluding it would conceal a real threat to ADB-based UI-test validity.

Future automation should select accessibility nodes or stable resource identifiers instead of relying on raw coordinates. It should record display size, application bounds, system-bar insets, orientation, density, and font scale for every run. Screenshots should be associated with the installed package hash and build hash so that visual evidence cannot drift from the tested binary.

\subsection{Visual Review Method}

Visual assessment considered hierarchy, identity, first-action clarity, state distinction, provenance visibility, limitation placement, empty states, populated states, navigation consistency, destructive-action treatment, and overall coherence. The assessment did not score aesthetic preference numerically because no validated participant scale was administered.

The overview succeeded by making the product proposition and review action dominant while keeping the evidence-reach warning in the first viewport. Summary cards used count, icon, label, and colour. Findings showed synthetic source and active pack near the decision. Settings gave the pack caveat and controlled demonstration their own sections. The custom navigation remained legible against the off-white background and used a clear selected state.

The result satisfies a project-defined candidate gate rather than a universal usability standard. One evaluator and one physical device cannot establish willingness to adopt across a population. The strongest supported statement is that the rendered candidate is coherent, presents its boundaries, and is suitable for internal-track and participant evaluation.

\subsection{Result Boundary}

The evidence supports the claim that the specified code compiles, the rule engine conforms to its fixtures, the Android artefacts build, the final package installs and launches, the main interface renders coherently, and provenance remains visible in the tested workflows. It does not support claims about 24-hour execution, battery cost, memory leaks, unrestricted third-party observation, multiple OEMs, legal compliance, accessibility conformance, production signing, or store acceptance.

\subsection{Evidence Triangulation}

No single method in the campaign was designed to carry the complete acceptance claim. Source inspection is good at confirming constants, branches, labels, and data flow, but it cannot establish which code entered an APK or how a screen rendered. Automated tests execute specified behaviours repeatedly, but their assertions can share assumptions with the implementation. Package inspection observes the assembled artefact, but it cannot prove the meaning of every runtime path. Device interaction observes a concrete installation, but the sample is narrow. The evaluation therefore uses convergence among methods while preserving the scope of each.

The deterministic tests and source review converge on threshold semantics. Source review identifies the pack-provided baseline, multiplier, daily threshold, burst limit, and source restrictions. Boundary assertions then test values immediately around those decision points. The fixed corpus tests varied inputs at scale. Agreement across these methods reduces the chance that the reported rule differs accidentally from the implemented one. It still does not establish that the rule is a valid legal classifier.

Runtime-boundary tests and presentation tests converge on uncertainty handling. Malformed evidence is rejected at admission, low-context evidence becomes an evidence gap, and presentation projection retains that state. This end-to-end check matters because a correct engine could be undermined if the UI translated an unknown state into reassuring copy. The tests show that selected fixtures preserve meaning through the current projection; screenshot review confirms the principal labels are visible in the rendered flow.

Build logs, output existence, checksums, manifest inspection, installed-package metadata, and physical launch converge on artefact identity. Each closes a different substitution risk. A successful Gradle message without an output file would be insufficient. A local file without a checksum could be confused after rebuilding. Source manifest facts could differ after dependency merging. An installed package could be a stale version. Reading these signals together supports the conclusion that the evaluated handset ran the recorded release candidate.

Visual review and UI-hierarchy inspection also complement each other. A screenshot shows layout, colour, clipping, and hierarchy as a person sees them. A hierarchy dump exposes text, bounds, and accessibility-node structure that may not be obvious visually. Neither replaces screen-reader interaction, large-font testing, or participant comprehension. Their combined use supports a bounded expert assessment and identifies the next accessibility evidence required.

Triangulation is not vote counting. If one authoritative check contradicts another, the discrepancy must be resolved rather than outvoted. For example, an installed-package dump showing a sensitive permission would invalidate a source-only conclusion even if several documents said otherwise. The evaluation gives greater weight to evidence closest to the released artefact for package claims and to independent human or legal review for substantive claims.

\subsection{Oracle Independence and Correlated Defect Risk}

The main corpus labels are derived from the same published threshold model the implementation is intended to realise. This is appropriate for conformance testing: it asks whether the program implements the specification. It is insufficient for empirical validation because a mistaken specification can yield perfect agreement. The measured 1.0000 precision and recall are therefore implementation-oracle agreement values, not estimates of performance on real legal cases.

Several design choices reduce accidental coupling. The generator records expected polarity directly from its construction, a fresh engine limits cross-case state, boundary cases are written explicitly, and extended tests target properties beyond the main daily-total rule. Nevertheless, developer knowledge shaped both sides. The campaign did not use a separately developed executable oracle, blinded labels, or an external benchmark.

Mutation testing would provide one useful next check. Deliberately changing a comparison from greater-than to greater-than-or-equal, weakening source validation, removing replay replacement, or allowing overlapping history should cause specific tests to fail. Surviving mutations would identify assertions that execute code without distinguishing correct behaviour. Property-based generation could then explore many valid and invalid shapes while preserving stated invariants.

An independent reference implementation could further reduce correlated defects. It should be written from the regulation-pack contract by a reviewer who has not read the production engine, ideally in a different language or style. Both implementations could evaluate a versioned interchange corpus. Disagreement would trigger case analysis rather than majority voting. This process would validate formal semantics, though it would still not validate legal relevance.

For legal or ecological validation, the oracle must change. Controlled event traces can provide ground truth about whether acquisition observed an action. Independent privacy reviewers can label whether a scenario warrants investigation, using a documented protocol and recording disagreement. Neither source should be described as perfect legal truth. The evaluation question, sampling frame, and annotation uncertainty must be reported alongside performance metrics.

This distinction explains why the thesis retains the term deterministic corpus rather than benchmark. A benchmark normally suggests comparison against an external reference population. The current corpus is a rigorous regression artefact that protects rule semantics as the code changes. Its value is high, but different from an independently labelled real-world set.

\subsection{Metric Interpretation and Error Analysis}

Precision and recall were included because they expose false-positive and false-negative cells more clearly than accuracy in an imbalanced corpus. The designed 80/20 mixture would allow a trivial always-positive strategy to appear 80 per cent accurate. The full confusion matrix prevents that reading: all 200 controls were negative and all 800 constructed signals were positive under the specified oracle.

The exact totals also make the experiment reproducible. The campaign did not sample a changing external population; it evaluated a fixed generated set. Confidence intervals around the observed proportions would describe repeated sampling from an assumed generator, not uncertainty about the fixed run. More importantly, they would not address specification validity. For this reason the report presents exact agreement and claim boundaries instead of using a narrow interval to imply real-world certainty.

Error analysis remains meaningful even when the main matrix has no disagreements. Adversarial suites reveal categories the primary generator does not represent: unsafe numeric values, prototype-derived identifiers, malformed context, source forgery, temporal overlap, replay, and history isolation. These tests show that corpus size alone is not coverage. A small, carefully targeted case can provide more security evidence than hundreds of additional values drawn from the same safe range.

Future field studies should report more than aggregate precision and recall. Results should be stratified by permission, application category, source authority, foreground status, OEM, Android version, and window resolution. Reviewers should inspect false positives and negatives for common causes. Missing acquisition records should not be treated as model negatives. Cases excluded because evidence is insufficient should be counted and explained rather than removed from the denominator without notice.

Calibration should be evaluated separately from discrimination. A status hierarchy may rank more concerning cases correctly while its displayed wording still conveys excessive certainty. Participant studies can test whether confidence changes appropriately between authoritative, synthetic, and incomplete evidence. This human calibration outcome is central to the product's trust claim and cannot be inferred from the engine matrix.

Operational metrics also need denominators. A launch success rate requires a defined number of clean starts and devices. A crash-free percentage requires a meaningful observation period and collection method. A jank percentage depends on a representative interaction trace and enough rendered frames. The current evaluation reports bounded events directly rather than producing population-style percentages from a short sample.

\subsection{Store-Candidate Evidence Review}

The project-defined store-candidate gate was assessed across product identity, core journey, disclosure, packaging, privacy posture, accessibility foundations, and operational ownership. Product identity was supported by a consistent name, icon, visual language, and package identifier. Core journey was supported by first-launch overview, findings navigation, detailed rationale, pack selection, demonstration control, and local-data deletion. Disclosure was supported by visible source labels, pack identity, evidence limits, and non-verdict language.

Packaging evidence included release APK and AAB generation, target SDK 36, ABI outputs, final checksums, manifest reconciliation, and installation of the APK. Privacy posture included no sensitive runtime permission in the final package, disabled backup, local-first storage, and a draft privacy policy consistent with the tested build. Accessibility foundations included textual state labels, non-colour cues, labelled controls, target sizing intent, and a simple navigation structure.

The gate was not interpreted as proof that every store prerequisite is complete. Production upload signing was absent. The privacy policy required owner-controlled HTTPS publication. Console declarations, listing assets, support details, content rating, internal-track delivery, and store review remained external. These gaps do not negate the engineering candidate; they determine why it must not be described as a publicly released production version.

The visual adjectives in the gate were operationalised cautiously. ``Beautiful'' meant coherent hierarchy, spacing, typography, colour restraint, branded identity, and absence of visible broken states in the inspected flows. ``Trustworthy'' meant that evidence provenance and limitations were prominent, destructive controls were explicit, and the product did not present a legal badge. ``Willing to use'' meant that an evaluator could understand the purpose, reach findings, inspect reasons, switch the review framework, and control data without encountering a blocking usability defect.

These definitions create a repeatable expert gate, but they do not convert subjective judgement into population evidence. A single evaluator can miss cultural preferences, accessibility barriers, unfamiliar terminology, or long-term fatigue. The appropriate next step is internal-track and participant evaluation using the actual distributed candidate. The expert gate answers whether that investment is justified, not whether the market has accepted the product.

The candidate judgement also depends on the declared product category. A local research and accountability aid can be useful with bounded imported or simulated evidence. A product marketed as universal device surveillance would fail because the acquisition authority is not established. Store copy must preserve the former positioning. Changes to marketing are therefore capable of invalidating readiness even if the binary remains identical.

\subsection{Evaluation Repeatability and Change Control}

The evidence applies to a specific source revision and generated artefacts. Repeating only the final UI flow after changing engine code would not renew rule evidence. Rebuilding only the APK after changing privacy copy would not renew screenshot evidence. The project therefore treats acceptance as a dependency graph: a change invalidates every downstream artefact that incorporates or presents it.

Engine or pack changes require TypeScript compilation, logical suites, boundary suites, presentation projections, release rebuilds, hashes, and targeted device review. Android manifest or Gradle changes require merged-manifest inspection, package dump reconciliation, rebuilds, installation, and launch. UI changes require lint, package build, hierarchy and screenshot capture, accessibility checks, and renewed expert review. Policy or store-description changes require reconciliation with package behaviour and data flows.

Reproducible scripts reduce omission but do not remove the need to inspect their outputs. A command may exit successfully while testing stale generated code, selecting the wrong device, or reading a prior file. The campaign therefore compiles generated runners immediately before execution, checks package identity after installation, and associates evidence with checksums. Future automation should emit a machine-readable manifest linking commit, tool versions, commands, artefact hashes, device facts, and screenshots.

Environment details also matter. Windows path length and native dependency layout affected release compilation. The chosen workaround and architecture flag belong in the build record because a different environment may produce different outputs. Tool warnings about future Gradle compatibility remain maintenance observations even when the current build passes.

A release tag should be immutable, while packs and policies use explicit version identifiers. If a later investigation finds an error, the project should publish a successor and a correction note rather than rewriting the historical evidence. This thesis follows the same principle by creating Revision 10 chapter files instead of replacing Revision 9. Preservation makes it possible to see which claims changed and why.

Repeatability finally includes honest failure recording. Network, dependency, device-coordinate, and build-path failures can be part of the evidence because they expose operational assumptions. Removing them from the narrative would make future reproduction harder. A trustworthy evaluation explains both the final pass and the conditions needed to obtain it.

\input{Iteration-V2/Chapter6/chapter-06-10}
\input{Iteration-V2/Chapter7/chapter-07-10}
\bibliographystyle{ieeetr}
\bibliography{reference}
\end{document}
